\documentclass[a4paper,11pt]{article}
\usepackage[utf8]{inputenc}
\usepackage[T1]{fontenc}
\usepackage{lmodern}
\usepackage[margin=2.5cm]{geometry}
\usepackage{amsmath,amssymb,amsthm}
\usepackage{booktabs}
\usepackage{xurl}
\usepackage[hidelinks]{hyperref}
\hypersetup{
  pdftitle={Integrate once, merge exactly: accumulator-based data reduction for mapping and time-resolved scattering experiments},
  pdfauthor={Jérôme Kieffer, Edgar Gutiérrez-Fernández, Manfred Burghammer, Jonathan Wright, Loïc Huder, Thomas Vincent (ESRF)},
  pdfsubject={Exact merge of pre-integrated scattering data with the azimuthal error model},
}

\newtheorem{proposition}{Proposition}
\newcommand{\mean}[1]{\widehat{#1}}
\newcommand{\Om}{\Omega}
\newcommand{\OmS}{\Omega^{(2)}}
\newcommand{\D}[2]{D^{#1}_{#2}}
\newcommand{\AB}{{A\sqcup B}}
\newcommand{\Ab}{{A\sqcup\{b\}}}

\title{Integrate once, merge exactly: accumulator-based data reduction for mapping and time-resolved scattering experiments}
\author{Jérôme Kieffer, Edgar Gutiérrez-Fernández, Manfred Burghammer,\\
Jonathan Wright, Loïc Huder and Thomas Vincent\\[0.5em]
\normalsize European Synchrotron Radiation Facility (ESRF), Grenoble, France}
\date{\today}

\begin{document}
\maketitle

\begin{abstract}
Schubert and Gertz~\cite{SG18} derive single-pass, numerically stable and mergeable formulas for
the weighted mean and (co-)variance, where the same weights $\omega_i$ are used for the mean and for
the sum of deviation products. The \texttt{azimuthal} error model of pyFAI weights the mean with
$\omega_i$ but the deviations with $\omega_i^2$. This note re-derives Section~3 of~\cite{SG18} in this
setting. A cross-term appears, which vanishes in~\cite{SG18}: an exact merge of two partitions
requires one additional accumulator per variable, $\D{V}{P} = \sum_{i\in P}\omega_i^2(v_i-\mean{v}_P)$.
Equation numbers of the form (SG$n$) refer to equation $n$ of~\cite{SG18}.
The motivation is to store pre-integrated accumulators instead of raw images for maps and time
series, which saves about one order of magnitude in storage. A simulation validates the approach, and
two experimental examples are presented: the decrease of the spottiness of a powder pattern with the
number of merged frames (ID11, ESRF), and the measurement of texture and orientation in a Kevlar fibre
(ID13, ESRF).
\end{abstract}

\section{Motivation: saving storage for maps and time series}

In a mapping or time-resolved experiment, all frames are first integrated. The map built from these
results is used to define regions with equivalent signal, and the frames of each region are then
combined to improve the statistics. Two workflows are possible:
\begin{enumerate}
\item The raw frames are kept. For each region, they are averaged and integrated a second time.
\item Each frame is integrated once, and the per-bin accumulators are stored in addition to
  $(q, I, \sigma)$. For each region, the pre-integrated accumulators of the selected frames are merged,
  without going back to the raw data.
\end{enumerate}

As an example, consider $10^6$ frames from an Eiger~4M detector ($2167\times2070$ pixels, 32-bit
integers), integrated in 2D on $4000$ radial $\times$ $8$ azimuthal bins. We assume a compression ratio of
8 for the raw integer data and of 3 for floating-point data. The accumulators per bin are the weighted
sum of signal $X_P$, the sum of weights $\Om_P$, the sum of squared weights $\OmS_P$, the pixel count,
the sum of squared deviations $XX_P$ and the new accumulator $\D{X}{P}$ introduced in
Section~\ref{sec:notation}.

\begin{table}[ht]
\centering
\small
\begin{tabular}{lrrr}
\toprule
Stored data & per frame & $10^6$ frames & gain \\
\midrule
Workflow 1: raw frames (32-bit) $+$ $I$, $\sigma$ (float32) & 2.33\,MB & 2.33\,TB & 1 \\
Workflow 2: 6 accumulators (float32) & 256\,kB & 256\,GB & 9.1 \\
Workflow 2: 6 accumulators (float64) & 512\,kB & 512\,GB & 4.5 \\
Workflow 2: $X_P$, $XX_P$, $\D{X}{P}$ (float32); $\Om_P$, $\OmS_P$, count shared & 128\,kB & 128\,GB & 18.2 \\
\bottomrule
\end{tabular}
\caption{Storage needed by both workflows for $10^6$ Eiger~4M frames integrated on $4000\times8$ bins.}
\label{tab:storage}
\end{table}

Table~\ref{tab:storage} shows that storing the accumulators saves an order of magnitude compared to
keeping the raw frames. Compared to storing only $(I,\sigma)$ (85\,GB), the overhead is 43--171\,GB,
i.e.\ 2--8\,\% of the raw data volume. $\Om_P$, $\OmS_P$ and the count only depend on the geometry, the
mask and the normalization (solid angle, polarization, flat-field). When those are static, they can be
stored once, a per-frame monitor being a scalar factor ($\Om_P\propto m$, $\OmS_P\propto m^2$). Dynamic
masks or sigma-clipping make them frame-dependent. The noisy accumulators $XX_P$ and $\D{X}{P}$ may also
compress less than assumed. The second pass is cheaper as well: about 9 times less data is read, and
the merge of $32\,000$ bins replaces the integration of $4.5$\,Mpixels per frame.

Workflow 2 requires an exact merge of the accumulators with the \texttt{azimuthal} error model, which
is the subject of this note.

\section{Why squared weights in pyFAI}

In pyFAI, each pixel $i$ provides a raw signal $s_i$ and a normalization $n_i$ (solid angle,
polarization, flat-field, \dots), and contributes to a bin with a pixel-splitting coefficient $c_i$.
With $x_i = s_i/n_i$ and $\omega_i = c_i n_i$, the integrated intensity of a bin $P$ is the weighted mean
\begin{equation}
  \mean{x}_P = \frac{\sum_{i\in P} c_i s_i}{\sum_{i\in P} c_i n_i}
             = \frac{\sum_{i\in P}\omega_i x_i}{\sum_{i\in P}\omega_i}.
\end{equation}
Neglecting the uncertainty of the denominator, its variance is
$\sum_{i\in P} c_i^2\,\sigma^2(s_i) \big/ \Om_P^2$.
The \texttt{poisson} error model propagates a known $\sigma^2(s_i)$, while the \texttt{azimuthal}
model estimates it from the dispersion within the bin: $\sigma^2(s_i)\approx n_i^2(x_i-\mean{x}_P)^2$.
Both models then share the same form:
\begin{equation}
  \sigma^2(\mean{x}_P) \approx \frac{1}{\Om_P^2}\sum_{i\in P}\omega_i^2 (x_i-\mean{x}_P)^2 = \frac{XX_P}{\Om_P^2},
\end{equation}
so that pyFAI reports $\mathrm{sem}=\sqrt{XX_P}/\Om_P$ and $\mathrm{std}=\sqrt{XX_P/\OmS_P}$.
The mean is weighted by $\omega_i$, the sum of squared deviations by $\omega_i^2$.

\section{Notation}
\label{sec:notation}

Table~\ref{tab:notation} extends Table~1 of~\cite{SG18}. The only changes are the sum of squared
weights $\OmS_P$, the $\omega^2$ weighting of the deviation products $VW_P$, and the new first-order
deviation sums $\D{V}{P}$.

\begin{table}[ht]
\centering
\small\setlength{\tabcolsep}{4pt}\begin{tabular}{lll}
\toprule
Symbol & Meaning & pyFAI ($V=W=X$) \\
\midrule
$A, B, P$ & partitions of the data set & bins, sub-bins \\
$\AB$ & union of disjoint partitions $A$ and $B$ & \\
$\Ab$ & partition $A$ plus a single point $b\notin A$ & \\
$X=\{x_i\}$, $Y$, $V$, $W$ & variables & $x_i = s_i/n_i$ \\
$\Om=\{\omega_i\}$ & weights & $\omega_i = c_i n_i$ \\
$V_P=\sum_{i\in P}\omega_i v_i$ & weighted sum of $V$ & \texttt{sum\_signal} \\
$\Om_P=\sum_{i\in P}\omega_i$ & sum of weights & \texttt{sum\_normalization} \\
$\OmS_P=\sum_{i\in P}\omega_i^2$ & sum of squared weights & \texttt{sum\_normalization2} \\
$\mean{v}_P=\frac{1}{\Om_P}V_P$ & weighted mean of $V$ & \texttt{intensity} \\
$VW_P=\sum_{i\in P}\omega_i^2(v_i-\mean{v}_P)(w_i-\mean{w}_P)$ & $\omega^2$-weighted sum of deviation products & \texttt{sum\_variance} \\
$\D{V}{P}=\sum_{i\in P}\omega_i^2(v_i-\mean{v}_P)$ & $\omega^2$-weighted sum of deviations & \emph{new accumulator} \\
\bottomrule
\end{tabular}
\caption{Notation, adapted from Table~1 of~\cite{SG18}.}
\label{tab:notation}
\end{table}

In~\cite{SG18}, $VW_P=\sum_{i\in P}\omega_i(v_i-\mean{v}_P)(w_i-\mean{w}_P)$, and the analogue of
$\D{V}{P}$, namely $\sum_{i\in P}\omega_i(v_i-\mean{v}_P)$, is identically zero by definition of the
mean. With squared weights,
\begin{equation}
  \D{V}{P} = \sum_{i\in P}\omega_i^2 v_i - \OmS_P\,\mean{v}_P = \OmS_P\left(\mean{v}^{(\omega^2)}_P - \mean{v}_P\right),
  \label{eq:D}
\end{equation}
where $\mean{v}^{(\omega^2)}_P$ is the $\omega^2$-weighted mean. $\D{V}{P}$ vanishes when both means
coincide, in particular when all weights within $P$ are equal, but not in general.

\section{Updateable weighted (co-)variance}

\subsection{Mean and ``textbook'' form}

The sums $V_P$, $\Om_P$ and $\OmS_P$ are additive over disjoint partitions, (SG14)--(SG15):
\begin{equation}
  V_\AB = V_A + V_B, \qquad \Om_\AB = \Om_A + \Om_B, \qquad \OmS_\AB = \OmS_A + \OmS_B.
\end{equation}
The update of the mean is unchanged, (SG16)--(SG17):
\begin{equation}
  \mean{v}_\AB = \frac{\Om_A}{\Om_\AB}\mean{v}_A + \frac{\Om_B}{\Om_\AB}\mean{v}_B
               = \mean{v}_A + \frac{\Om_B}{\Om_\AB}\left(\mean{v}_B - \mean{v}_A\right).
  \label{eq:mean}
\end{equation}
Expanding the product as in (SG18)--(SG19) gives the analogue of the König--Huygens formula:
\begin{equation}
  VW_P = \sum_{i\in P}\omega_i^2 v_i w_i - \mean{w}_P\sum_{i\in P}\omega_i^2 v_i
         - \mean{v}_P\sum_{i\in P}\omega_i^2 w_i + \OmS_P\,\mean{v}_P\,\mean{w}_P .
  \label{eq:textbook}
\end{equation}
Unlike (SG19), the last three terms do not collapse into a single product, as $\mean{v}_P$ is not the
$\omega^2$-weighted mean. Like the ``textbook'' algorithm, Equation~\eqref{eq:textbook} only involves
plain sums, which are trivially mergeable, but it suffers from catastrophic cancellation when
$VW_P \ll \OmS_P\,\mean{v}_P\,\mean{w}_P$. It should not be used.

\subsection{Derivation of the general form}
\label{sec:derivation}

Section~3.1 of~\cite{SG18} proves (SG21) through an identity derived ``backwards''. With squared
weights, a direct decomposition is simpler. Let $\delta^v_A = \mean{v}_A - \mean{v}_\AB$ be the shift of
the mean of $A$ with respect to the mean of the union. Writing
$v_i-\mean{v}_\AB = (v_i-\mean{v}_A)+\delta^v_A$ and similarly for $w$:
\begin{align}
  \sum_{i\in A}\omega_i^2(v_i-\mean{v}_\AB)(w_i-\mean{w}_\AB)
  &= \sum_{i\in A}\omega_i^2\left[(v_i-\mean{v}_A)+\delta^v_A\right]\left[(w_i-\mean{w}_A)+\delta^w_A\right]\nonumber\\
  &= VW_A + \delta^w_A\,\D{V}{A} + \delta^v_A\,\D{W}{A} + \OmS_A\,\delta^v_A\,\delta^w_A .
\end{align}
In~\cite{SG18}, the two middle terms vanish. Adding the same expression for $B$ gives the
``parallel axis'' form:
\begin{equation}
  VW_\AB = VW_A + VW_B + \sum_{P\in\{A,B\}}\left(\OmS_P\,\delta^v_P\,\delta^w_P + \delta^w_P\,\D{V}{P} + \delta^v_P\,\D{W}{P}\right).
  \label{eq:parallel}
\end{equation}
From Equation~\eqref{eq:mean}, with $\Delta v = \mean{v}_A - \mean{v}_B$:
\begin{equation}
  \delta^v_A = \frac{\Om_B}{\Om_\AB}\,\Delta v, \qquad
  \delta^v_B = -\frac{\Om_A}{\Om_\AB}\,\Delta v .
\end{equation}
The same decomposition applied to $\D{V}{\AB}=\sum_{i\in\AB}\omega_i^2(v_i-\mean{v}_\AB)$ gives
$\D{V}{\AB} = \D{V}{A}+\D{V}{B}+\OmS_A\,\delta^v_A+\OmS_B\,\delta^v_B$, hence:

\begin{proposition}[General form, analogue of (SG21)]
For disjoint partitions $A$ and $B$, with $\Delta v = \mean{v}_A-\mean{v}_B$ and $\Delta w = \mean{w}_A-\mean{w}_B$:
\begin{align}
  VW_\AB &= VW_A + VW_B
           + \frac{\OmS_A\,\Om_B^2 + \OmS_B\,\Om_A^2}{\Om_\AB^2}\,\Delta v\,\Delta w \nonumber\\
         &\quad + \frac{\left(\Om_B\D{V}{A}-\Om_A\D{V}{B}\right)\Delta w + \left(\Om_B\D{W}{A}-\Om_A\D{W}{B}\right)\Delta v}{\Om_\AB}, \label{eq:general}\\
  \D{V}{\AB} &= \D{V}{A} + \D{V}{B} + \frac{\OmS_A\,\Om_B - \OmS_B\,\Om_A}{\Om_\AB}\,\Delta v. \label{eq:Dmerge}
\end{align}
\end{proposition}

\paragraph{Variance.} For $V=W=X$, the case used by pyFAI:
\begin{align}
  XX_\AB &= XX_A + XX_B + \frac{\OmS_A\,\Om_B^2 + \OmS_B\,\Om_A^2}{\Om_\AB^2}\,\Delta x^2
           + 2\,\frac{\Om_B\D{X}{A}-\Om_A\D{X}{B}}{\Om_\AB}\,\Delta x, \label{eq:variance}\\
  \D{X}{\AB} &= \D{X}{A} + \D{X}{B} + \frac{\OmS_A\,\Om_B - \OmS_B\,\Om_A}{\Om_\AB}\,\Delta x . \label{eq:DXmerge}
\end{align}
Equation~\eqref{eq:variance} is symmetric in $A$ and $B$: unlike the current pyFAI implementation, the
partitions do not need to be ordered by weight.

\paragraph{Sufficient statistics.}
In~\cite{SG18}, the four values $(VW_P, \mean{v}_P, \mean{w}_P, \Om_P)$ are sufficient statistics for
merging. With squared weights, $(VW_P, \mean{v}_P, \mean{w}_P, \Om_P, \OmS_P, \D{V}{P}, \D{W}{P})$ are
needed. For the variance alone, this means five accumulators per bin: $X_P$, $\Om_P$, $\OmS_P$, $XX_P$ and
$\D{X}{P}$. $\D{X}{P}$ cannot be dropped: two partitions may share the same
$(X_P, \Om_P, \OmS_P, XX_P)$ but have different $\D{X}{P}$, and therefore give different $XX$ once merged
with the same third partition. No formula based on these four accumulators can be exact.

\paragraph{Reduction to~\cite{SG18}.}
If $\D{V}{A}=\D{V}{B}=\D{W}{A}=\D{W}{B}=0$ and $\OmS$ is replaced by $\Om$ (same weights for mean and
deviations), the coefficient of $\Delta v\,\Delta w$ becomes
$(\Om_A\Om_B^2+\Om_B\Om_A^2)/\Om_\AB^2 = \Om_A\Om_B/\Om_\AB$, and Equation~\eqref{eq:general} reduces to (SG21).

\paragraph{Form using sums.}
As in (SG22), the means can be replaced by the sums, using
$\Delta v = (\Om_B V_A - \Om_A V_B)/(\Om_A\Om_B)$:
\begin{multline}
  VW_\AB = VW_A + VW_B
   + \frac{\left(\OmS_A\Om_B^2+\OmS_B\Om_A^2\right)\left(\Om_B V_A-\Om_A V_B\right)\left(\Om_B W_A-\Om_A W_B\right)}{\Om_A^2\,\Om_B^2\,\Om_\AB^2}\\
   + \frac{\left(\Om_B\D{V}{A}-\Om_A\D{V}{B}\right)\left(\Om_B W_A-\Om_A W_B\right)
          + \left(\Om_B\D{W}{A}-\Om_A\D{W}{B}\right)\left(\Om_B V_A-\Om_A V_B\right)}{\Om_A\,\Om_B\,\Om_\AB}.
  \label{eq:sums}
\end{multline}
This is the form to use with the accumulators of pyFAI, which stores sums rather than means.

\subsection{Specializations to a single point}

Section~3.2 of~\cite{SG18} specializes the general form to a partition $B=\{b\}$ made of a single point.
Its statistics are $V_b=\omega_b v_b$, $\Om_b=\omega_b$, $\OmS_b=\omega_b^2$, $VW_b=0$ and also
$\D{V}{b}=0$. With $\Delta v = \mean{v}_A - v_b$:
\begin{align}
  VW_\Ab &= VW_A + \omega_b^2\,\frac{\OmS_A + \Om_A^2}{\Om_\Ab^2}\,\Delta v\,\Delta w
            + \omega_b\,\frac{\D{V}{A}\,\Delta w + \D{W}{A}\,\Delta v}{\Om_\Ab}, \label{eq:single}\\
  \D{V}{\Ab} &= \D{V}{A} + \omega_b\,\frac{\OmS_A - \omega_b\,\Om_A}{\Om_\Ab}\,\Delta v .
\end{align}
The term in $\D{V}{A}$ remains: even an online algorithm, which adds one point at a time, has to carry
$\D{V}{A}$. Only the contribution of the new point vanishes.

\paragraph{Analogue of Welford's update (SG27).} Welford~\cite{Welford62}, in the form given by Knuth, updates the sum of squared deviations with the product of the deviations to the old and new means.
For $V=W=X$, using $\mean{x}_\Ab - x_b = \frac{\Om_A}{\Om_\Ab}(\mean{x}_A - x_b)$, Equation~\eqref{eq:single} reads
\begin{equation}
  XX_\Ab = XX_A + \omega_b^2\,\frac{\OmS_A + \Om_A^2}{\Om_A\,\Om_\Ab}\left(\mean{x}_\Ab - x_b\right)\left(\mean{x}_A - x_b\right)
           + 2\,\omega_b\,\frac{\D{X}{A}}{\Om_\Ab}\left(\mean{x}_A - x_b\right).
  \label{eq:welford}
\end{equation}
The update currently implemented in pyFAI is (SG27) with $\omega_b^2$ in place of $\omega_b$:
$XX_\Ab \approx XX_A + \omega_b^2(\mean{x}_\Ab - x_b)(\mean{x}_A - x_b)$. Compared to
Equation~\eqref{eq:welford}, it misses the factor $(\OmS_A+\Om_A^2)/(\Om_A\Om_\Ab)$ and the $\D{X}{A}$ term.
Both corrections vanish when all the weights of $A\sqcup\{b\}$ are equal: $\OmS_A = \omega_b\Om_A$ and
$\D{X}{A}=0$. This is why the current formula is accurate for nearly uniform weights, and in particular
for unweighted data.

\paragraph{Analogue of Youngs and Cramer's update (SG28).} Youngs and Cramer~\cite{YC71} use the running sum instead of the running mean.
Using the sum $X_A$ instead of the mean, with $\mean{x}_A - x_b = (X_A - \Om_A x_b)/\Om_A$:
\begin{equation}
  XX_\Ab = XX_A + \omega_b^2\,\frac{\OmS_A + \Om_A^2}{\Om_A^2\,\Om_\Ab^2}\left(X_A - \Om_A x_b\right)^2
           + 2\,\omega_b\,\frac{\D{X}{A}}{\Om_A\,\Om_\Ab}\left(X_A - \Om_A x_b\right).
  \label{eq:youngs}
\end{equation}
As in~\cite{SG18}, the division comes last and only $XX$ and $\D{X}{}$ depend on it. This is the
preferred form for pipelined CPUs and GPUs.

\subsection{Numerical remarks}

\begin{itemize}
\item $\D{V}{P}$ could be replaced by the plain sum $\sum_{i\in P}\omega_i^2 v_i$, which is trivially
  mergeable (Equation~\eqref{eq:D}). However, recovering $\D{V}{P}$ from it involves the difference of two
  large, close numbers. The centered quantity $\D{V}{P}$ is preferred for the same reason as $VW_P$ is
  preferred to $\sum\omega_i^2 v_i w_i$.
\item $XX_P$ and $\OmS_P$ are sums of positive terms and are insensitive to cancellation: single
  precision is sufficient for an uncertainty. $\D{X}{P}$ is a correction term, also accurate enough in
  single precision. Conversely, $\Delta x = \mean{x}_A - \mean{x}_B$ is a difference of close numbers:
  $X_P$ and $\Om_P$ deserve the highest available precision (double-word arithmetic on GPU).
\item In the OpenCL implementation, the \texttt{float8} accumulator stores $X_P$, $XX_P$ and $\Om_P$ in
  double-word precision, the count and $\OmS_P$ in single precision. In the \texttt{azimuthal} model, the
  low word of $XX_P$ is the natural slot for $\D{X}{P}$.
\end{itemize}

\subsection{Scope of the exact merge: pixel splitting}
\label{sec:splitting}

Equations~\eqref{eq:variance} and~\eqref{eq:DXmerge} are exact for the union of the contributions of
$A$ and $B$. They reproduce a direct integration of $A\cup B$ only if no pixel contributes to both
partitions. This is always the case when merging different frames or scan points, since each (pixel,
frame) pair belongs to a single frame. It is not the case when merging neighbouring bins of the same
frame (azimuthal sectors into a 1D profile, 2D into 1D rebinning) if pixel splitting is enabled.

A pixel $i$ split between $A$ and $B$ with the coefficients $c_{iA}$ and $c_{iB}$ appears in the merge as
two contributions with the same $x_i$. $\Om_P$ and $X_P$ are linear in $c$ and remain exact. But $\OmS_P$,
$XX_P$ and $\D{X}{P}$ receive $(c_{iA}^2+c_{iB}^2)\,n_i^2$ instead of $(c_{iA}+c_{iB})^2\,n_i^2$. As
$c_{iA}^2+c_{iB}^2 < (c_{iA}+c_{iB})^2$, they are underestimated, and so is the propagated uncertainty.
Moreover, pixel splitting correlates neighbouring bins~\cite{Yang14}, which biases the $\chi^2$ of
Equation~\eqref{eq:chi2} when the split pixels are shared between azimuthal sectors.

\textbf{The merge is therefore rigorously exact only when pixel splitting is disabled} (pyFAI method
\texttt{("no", "csr", \dots)}), or when the merged partitions do not share any pixel. With hybrid
photon-counting detectors such as Eiger or Pilatus, whose point-spread function is about one pixel
with no correlation between neighbouring pixels, pixel splitting brings little and is best disabled for
data meant to be merged, with bins at least as wide as a pixel. With splitting enabled, the error remains
small when the bins are much wider than a pixel, since only the pixels on the bin edges are shared.

\section{Consequences for pyFAI}
\sloppy

\begin{itemize}
\item Every engine implementing the \texttt{azimuthal} error model (Cython CSR, CSC, LUT, histogram;
  OpenCL CSR) should accumulate $\D{X}{P}$ and use Equations~\eqref{eq:variance}, \eqref{eq:sums}
  or~\eqref{eq:youngs}, which are exact.
\item Integrated results should expose $\D{X}{P}$ next to $XX_P$ (\texttt{sum\_variance}) and $\OmS_P$
  (\texttt{sum\_normalization2}). Already-integrated results can then be merged exactly with the
  \texttt{azimuthal} error model: union of results, multi-geometry, and rebinning of 2D results into 1D
  (the latter only without pixel splitting, Section~\ref{sec:splitting}).
\item A numerical validation is provided in the tutorial notebook
  \nolinkurl{doc/source/usage/tutorial/Variance/uncertainties.ipynb}.
\end{itemize}

\section{Validation by simulation}
\label{sec:simulation}

\subsection{Setup}

Both workflows of Section~1 are compared on a single bin of $1000$ pixels, with a normalization
$n_i$ uniformly distributed in $[0.7, 1.3]$ and $\omega_i = n_i$. A series is made of $N=32$ independent
frames. Each frame has a background of $20$ counts and very ``spotty'' diffraction: each pixel has a
1\,\% probability to receive a spot of $2000$ counts, at positions drawn again for every frame. Poisson
noise is added to the expected signal $n_i(\text{background}+\text{spots})$.
\begin{itemize}
\item Workflow 1: the raw frames are averaged, then integrated once with the \texttt{azimuthal} model.
\item Workflow 2: the union of all (pixel, frame) pairs is integrated, which is exactly what the merge
  of the per-frame accumulators gives.
\end{itemize}
The experiment is repeated $1000$ times. The true standard error of the mean is the dispersion of the
integrated intensity over the repetitions; with $1000$ repetitions it is known to about $\pm2\,\%$.

Applying the decomposition of Section~\ref{sec:derivation} to the $N$ values of each pixel gives
\begin{equation}
  XX_{[2]} = N\,XX_{[1]} + \sum_{i}\omega_i^2\sum_{k=1}^{N}\left(x_{ik}-\bar{x}_i\right)^2,
  \label{eq:temporal}
\end{equation}
where $\bar{x}_i$ is the average of pixel $i$ over the frames. Workflow 1 removes the temporal dispersion
of each pixel (the last term) before integration, workflow 2 keeps it. With $\Om_{[2]}=N\,\Om_{[1]}$ (indices $[1]$ and $[2]$ refer to the workflows),
both give the same intensity, but their standard errors differ in general. Writing each pixel value as
a true value, plus a systematic azimuthal structure of variance $S$ (texture, constant over the frames),
plus a noise of variance $\sigma^2$ independent from frame to frame, one expects
\begin{equation}
  \mathrm{sem}_1^2 \approx \frac{\OmS}{\Om^2}\left(\frac{\sigma^2}{N} + S\right), \qquad
  \mathrm{sem}_2^2 \approx \frac{\OmS}{\Om^2}\left(\frac{\sigma^2}{N} + \frac{S}{N}\right).
\end{equation}
Without texture ($S=0$), both are equivalent. Spots appearing at random positions in every frame belong
to $\sigma^2$, not to $S$.

\subsection{Results}

Table~\ref{tab:simulation} gives the results. Equation~\eqref{eq:temporal} is verified to $2\times10^{-16}$,
and the intensities of both workflows agree to $7\times10^{-15}$.

\begin{table}[ht]
\centering
\small
\begin{tabular}{lcc}
\toprule
 & No texture & Fixed texture \\
\midrule
True standard error of the mean & 1.150 & 1.182 \\
Workflow 1 (average frames, then integrate) & $1.128 \pm 0.041$ & $1.211 \pm 0.043$ \\
Workflow 2 (merge accumulators) & $1.129 \pm 0.034$ & $1.131 \pm 0.032$ \\
\bottomrule
\end{tabular}
\caption{Standard error of the mean estimated by both workflows ($N=32$ frames, $1000$ pixels, $1000$
repetitions). The $\pm$ is the dispersion of the estimator over the repetitions.}
\label{tab:simulation}
\end{table}

\paragraph{Without texture.} Both workflows give the same standard error, in agreement with the true
value: both are about 2\,\% low, at the limit of the precision of the simulation, which may come from
the heavy-tailed distribution of the spotty signal. Workflow 2 is even slightly more stable
($\pm0.034$ instead of $\pm0.041$), as the dispersion is estimated from $32\,000$ samples instead of $1000$.

\paragraph{With a fixed texture.} To show the difference, a systematic structure is added: a Gaussian
random offset of standard deviation $15$ counts per pixel, drawn once per repetition and kept for its
$N$ frames. As expected, workflow 1 keeps the term $S$ and slightly overestimates the error ($1.211$),
whereas workflow 2 divides it by $N$ and underestimates it ($1.131$). It counts $N$ measurements of the
same pixels as independent samples.

\subsection{Checking for texture before merging}

The absence of texture can be checked on the pre-integrated data themselves, frame by frame, using the
azimuthal bins. Without texture, the intensities $I_j$ of the $M$ azimuthal sectors of a radial bin agree
within their standard errors $\mathrm{sem}_j$, and
\begin{equation}
  \chi^2 = \sum_{j=1}^{M}\frac{\left(I_j - I\right)^2}{\mathrm{sem}_j^2}
  \label{eq:chi2}
\end{equation}
is close to its expectation of $M-1$ (7 for 8 sectors), where $I$ is the intensity of the merged sectors. The same test can
be applied to the merge of several frames, where a systematic texture accumulates whereas random
spots average out. Once the absence of texture is validated, the azimuthal bins are merged with
Equations~\eqref{eq:variance} and~\eqref{eq:DXmerge}, which gives a single 1D integrated signal with an
exact propagation of the \texttt{azimuthal} error, provided pixel splitting is disabled
(Section~\ref{sec:splitting}). This test only detects a texture at the scale of the
azimuthal sectors. A texture at the scale of the pixels, as in the simulation above, cannot be seen in
the accumulators; it requires the raw frames, e.g.\ through the ratio of both terms of
Equation~\eqref{eq:temporal}.

\section{Experimental examples}

\subsection{Spottiness: is the sample exposed enough?}
\label{sec:spottiness}

Rietveld refinement assumes an ideal powder: each Debye--Scherrer ring is sampled by a large number of
randomly oriented crystallites, and its intensity is uniform along the azimuth. With a small beam, large
crystallites or a short exposure, rings break into spots and the integrated intensities become
unreliable. The \texttt{azimuthal} error model directly measures this: the relative standard error of
the mean of radial bin $r$ is
\begin{equation}
  \frac{\mathrm{sem}_r}{I_r} = \frac{\sqrt{XX_r}/\Om_r}{X_r/\Om_r} = \frac{\sqrt{XX_r}}{X_r},
\end{equation}
and pyFAI defines the spottiness index (\texttt{IntegrateResult.calc\_spottiness}), inspired
by~\cite{Ingham14}, as its quadratic mean over the radial bins:
\begin{equation}
  S = \sqrt{\frac{\sum_r w_r\,XX_r/X_r^2}{\sum_r w_r}},
  \label{eq:spottiness}
\end{equation}
with $w_r=1$ (unweighted) or $w_r=I_r$ (weighted by the intensity of the rings, which emphasizes the
Bragg peaks). As a rule of thumb, $S<0.05$ corresponds to a smooth powder pattern, $0.05<S<0.15$ to a
mild spottiness or texture, and $S>0.15$ to a strongly spotty pattern, likely with large grains.
$S$ only depends on the accumulators, so it can be calculated after any merge of pre-integrated results.

The dataset used here was recorded on the ID11 beamline of the ESRF on the standard reference material
SRM~640 (silicon), with a sub-micron beam at 40\,keV and an Eiger2~CdTe~4M detector at 200\,Hz. The first
data file contains 200 frames, out of 90 files. It is processed in the pyFAI tutorial
\nolinkurl{doc/source/usage/tutorial/Spotty.ipynb}, with $2000$ bins in $2\theta$ and the
\texttt{azimuthal} error model:
\begin{itemize}
\item a single frame is strongly spotty: $S\approx0.16$ (unweighted) and $S\approx0.21$ (weighted), with
  large fluctuations from frame to frame for the weighted index;
\item merging the pre-integrated results of successive frames decreases $S$ to about $0.02$ after 200
  frames. Beyond a few tens of frames, $S$ decreases as $1/\sqrt{N}$, as expected for independent frames
  (Section~\ref{sec:simulation}). The 199 merges take about $0.1$\,s;
\item integrating the average of the 200 raw frames (workflow 1) gives $S=0.025$ unweighted but
  $S=0.061$ weighted. According to Section~\ref{sec:simulation}, this difference points to a systematic
  azimuthal structure on the most intense rings, which survives the averaging. Its origin (residual
  texture, dynamically masked pixels of the detector, geometry) remains to be investigated.
\end{itemize}

The spottiness index is therefore a metric to decide when a sample has been exposed long enough:
frames are integrated on the fly, their accumulators merged, and the acquisition stops once $S$ goes
below a target value, e.g.\ $0.05$ for a Rietveld-quality pattern. For this dataset, about 35 frames are
needed, whereas a pure $1/\sqrt{N}$ extrapolation from the first frame would predict 10: the measured
curve, not the extrapolation, should be used. Exposing no longer than needed limits the radiation
damage to the sample and the cost of the experiment (beam time, storage). The tutorial still uses the
approximate merge formula currently implemented in pyFAI; the figures will be updated with the exact
merge of Section~\ref{sec:derivation}.

\subsection{Texture and orientation: a Kevlar fibre}
\label{sec:kevlar}

Texture is the opposite case: the azimuthal distribution of the intensity is the signal of interest,
and the azimuthal sectors must not be merged. The dataset, recorded on the ID13
beamline of the ESRF, is a tomographic scan of a Kevlar fibre, i.e.\ poly(\emph{p}-phenylene
terephthalamide), a highly oriented polymer~\cite{Northolt74}. Its main characteristics, read from the
data files, are:
\begin{itemize}
\item Eiger~4M detector (silicon sensor, $450\,\mu$m), $2167\times2070$ pixels, 16-bit frames;
\item photon energy $13.45$\,keV ($\lambda=0.922$\,\AA), $10$\,ms per frame;
\item $85\times160 = 13\,600$ frames, stored in 7 files of about $13.3$\,GB in total (compressed);
\item a calibration dataset on corundum ($\alpha$-Al$_2$O$_3$), with the corresponding calibration notebook.
\end{itemize}
The data are available at \url{http://ftp.edna-site.org/pyFAI/kevlar/}. This well-known dataset has not
been processed yet with the method of this note; the following describes the planned analysis.

\paragraph{Integration and storage.} Each frame is integrated on $1000$ radial $\times$ $36$ azimuthal
bins (sectors of $10^\circ$), with the \texttt{azimuthal} error model and without pixel splitting, so
that the sectors can later be merged exactly (Section~\ref{sec:splitting}), and the accumulators are stored.
With the assumptions of Section~1 (compression ratio of 3 for floating-point data), the 6 accumulators
in float32 take about $288$\,kB per frame, i.e.\ $3.9$\,GB for the whole scan, or $2.0$\,GB when
$\Om_P$, $\OmS_P$ and the count are shared. Compared to the $13.3$\,GB of raw data, the gain is only 3.4 to
6.8, smaller than in Section~1: the raw frames are 16-bit, and 36 azimuthal sectors are kept.

\paragraph{Texture.} For each frame and each radial bin, the $\chi^2$ of Equation~\eqref{eq:chi2} with
$M=36$ sectors (expectation $35$) tests the isotropy of the signal. Frames where the beam misses the
fibre (air scattering only) should be consistent with $\chi^2\approx35$, while frames through the fibre
should be strongly anisotropic. This validates the texture test on real data, and gives a map of the
anisotropy over the scan.

\paragraph{Orientation.} For the textured frames, the azimuthal profile $I_j$ of the strong equatorial
reflections of Kevlar, (110) and (200) at $d\approx4.3$ and $3.9$\,\AA~\cite{Northolt74}, gives the
orientation of the fibre axis projected on the detector (position of the maxima) and the degree of
orientation (width of the peaks, or Hermans orientation factor). The uncertainties $\mathrm{sem}_j$
propagated by the \texttt{azimuthal} model give error bars on these parameters. Frames with the same
orientation can then be merged sector by sector, exactly, to improve the statistics, without going back
to the raw data.

{\sloppy\begin{thebibliography}{9}
\bibitem{SG18}
E.~Schubert and M.~Gertz.
\newblock Numerically stable parallel computation of (co-)variance.
\newblock In \emph{Proceedings of the 30th International Conference on Scientific and Statistical
  Database Management (SSDBM~'18)}, pages 10:1--10:12, Bozen-Bolzano, Italy, 2018.
\newblock \href{https://doi.org/10.1145/3221269.3223036}{doi:10.1145/3221269.3223036}.

\bibitem{Welford62}
B.~P. Welford.
\newblock Note on a method for calculating corrected sums of squares and products.
\newblock \emph{Technometrics}, 4(3):419--420, 1962.
\newblock \href{https://doi.org/10.2307/1266577}{doi:10.2307/1266577}.

\bibitem{YC71}
E.~A. Youngs and E.~M. Cramer.
\newblock Some results relevant to choice of sum and sum-of-product algorithms.
\newblock \emph{Technometrics}, 13(3):657--665, 1971.
\newblock \href{https://doi.org/10.1080/00401706.1971.10488826}{doi:10.1080/00401706.1971.10488826}.
\bibitem{Ingham14}
B.~Ingham.
\newblock Statistical measures of spottiness in diffraction rings.
\newblock \emph{Journal of Applied Crystallography}, 47(1):166--172, 2014.
\newblock \href{https://doi.org/10.1107/S1600576713029713}{doi:10.1107/S1600576713029713}.

\bibitem{Yang14}
X.~Yang, P.~Juh\'as and S.~J.~L. Billinge.
\newblock On the estimation of statistical uncertainties on powder diffraction and small-angle
  scattering data from two-dimensional X-ray detectors.
\newblock \emph{Journal of Applied Crystallography}, 47(4):1273--1283, 2014.
\newblock \href{https://doi.org/10.1107/S1600576714010516}{doi:10.1107/S1600576714010516}.

\bibitem{Northolt74}
M.~G. Northolt.
\newblock X-ray diffraction study of poly(\emph{p}-phenylene terephthalamide) fibres.
\newblock \emph{European Polymer Journal}, 10(9):799--804, 1974.
\newblock \href{https://doi.org/10.1016/0014-3057(74)90131-1}{doi:10.1016/0014-3057(74)90131-1}.
\end{thebibliography}}

\end{document}
